\chapter{Conclusions}
\label{Conclusions}

Future work

 In this thesis work we have obtained two original results:
 1. ...
 2. ...
 These results are shown in Chapter ? and in Chapter ?, respectively. In particular,
 in Section ? we have shown... . Instead, in Sections ?
 and ? we have studied... 


 mencionar aixó: Several idealizations in the pipeline described above should be kept in mind when interpreting the results of this thesis. The noise realizations added to each injection are stationary and Gaussian (\texttt{gaussian-noise=True}), so none of the non-Gaussian transients \ref{sec:noise} (glitches) that affect real LIGO data are present; a real search would need to contend with this additional source of PE bias and false coincidences. The same projected O5a PSD is used for both H1 and L1, whereas the two real detectors have distinct, time-varying noise budgets; this symmetry is not expected to qualitatively change the sky localization results, which depend primarily on the relative antenna pattern and time delay between the two sites, but it does idealize the relative sensitivity assumed for each detector. Finally, the analysis assumes perfect detector calibration and does not marginalize over calibration uncertainty, an additional source of systematic error present in real LVK analyses that is not modelled here.
